\documentclass[12pt,a4paper,oneside]{article}
\usepackage[brazil]{babel}
\usepackage[utf8]{inputenc}
\usepackage{olymp}

\setcounter{problem}{4}

\begin{document}

\begin{problem}{Campo minado}{campo}{2}

Campo Minado é um popular jogo de computador para um único jogador. Ele é jogado em um tabuleiro $M \times N$ onde são escondidas diversas bombas, e o objetivo do jogador é descobrir onde estão escondidas todas as bombas. Ao clicar em uma casa do tabuleiro, o jogo revela para o jogador se aquela casa contém ou não uma bomba. Caso o jogador clique em uma casa com uma bomba, ele perde o jogo. 

Para facilitar, todas as casas já reveladas pelo jogador contêm um número, que indica quantas bombas existem nas oito (ou menos, para as casas na borda do tabuleiro) casas vizinhas a esta casa. Veja a imagem abaixo, onde as posições das bombas são indicadas por bandeiras. Casas vazias indicam que não há nenhuma bomba em suas casas vizinhas.

\begin{center}
\includegraphics[scale=0.4]{images/exercise_37_0.png}
\end{center}

Faça um programa que, dado o tamanho do tabuleiro e as posições das bombas, mostre a configuração final do jogo supondo que o jogador encontrou todas as bombas e revelou todas as demais casas do tabuleiro.

%\Notes
%
%Observações, se tiver.

\InputFile

A entrada contém diversas linhas. A primeira linha contém três inteiros $M$, $N$ e $B$, indicando respectivamente o número de linhas e colunas do tabuleiro, e o número de bombas escondidas. Em seguida virão $B$ linhas, cada uma com dois inteiros $Y_i$ e $X_i$, que indicam que a $i$-ésima bomba está escondida na casa que fica na linha $Y_i$ e na coluna $X_i$.

Restrições:
\begin{itemize}
    \item $1 \leq M , N \leq 100$
    \item $0\leq B \leq M \times N$
    \item $1 \leq Y_i \leq M$
    \item $1 \leq X_i \leq N$
\end{itemize}

\OutputFile

Seu programa deve imprimir $M$ linhas com $N$ caracteres cada, cada caractere representando uma das casas do tabuleiro. Para casas com bombas, imprima um `\texttt{B}'. Para casas em que não há nenhuma bomba e também não há nenhuma bomba em suas casas vizinhas, imprima um `\texttt{-}'. Para as demais casas, imprima um número indicando o número de bombas escondidas nas casas vizinhas.


% \Notes
% Preste atenção aos tipos das variáveis, pois $F(90) \approx 2^{61}$.

\Example

\begin{example}
\exmp{
9 9 10
1 7
2 3
2 8
4 1
4 2
5 4
6 6
6 7
6 9
9 7
}{
-111-1B21
-1B1-12B1
2321--111
BB211----
222B22221
--112BB2B
----12221
-----111-
-----1B1-
}%
\end{example}

%Comentários sobre o exemplo, se necessário.
Este exemplo corresponde ao da imagem mostrada acima.
\end{problem}

\end{document}
